%=============================================================================
% PAPER CLAIMS DOCUMENTARY RECORD
% Purely documentary: what section_alpha_locked.tex, appendix_K.tex, and the
% standalone PDF actually SAY, with line numbers and equation labels.
%
% This document does NOT attempt to reproduce any computation.
% It answers 10 specific questions about the paper's claims.
%
% Created: 2026-03-22
%=============================================================================
\documentclass[12pt,a4paper]{article}
\usepackage[margin=2.5cm]{geometry}
\usepackage{amsmath,amssymb}
\usepackage{booktabs,array,longtable}
\usepackage{hyperref}
\usepackage{tcolorbox}

\title{Paper Claims Documentary Record:\\[6pt]
\large What the DFD Alpha Derivation Actually States\\
(Line-by-Line, No Reproduction Attempted)}

\author{Documentary Audit}
\date{March 22, 2026}

\begin{document}
\maketitle

\begin{abstract}
This document provides a precise, line-by-line record of what the DFD
papers claim regarding the derivation of $\alpha^{-1} = 137.036$.
Three source files were read completely:
\texttt{section\_alpha\_locked.tex} (274 lines),
\texttt{appendix\_K.tex} (lines 1--563+ relevant to alpha),
and the standalone PDF ``Ab Initio Derivation of the Fine Structure
Constant from Density Field Dynamics'' (27 pages).
No computation is reproduced; this is purely documentary.
\end{abstract}

\tableofcontents
\newpage

%=============================================================================
\section{Question 1: Which Operator?}
\label{sec:q1}
%=============================================================================

\begin{tcolorbox}[colback=yellow!5!white, colframe=yellow!50!black,
  title=Answer: CONNECTION LAPLACIAN (scalar-type)]

\textbf{Source:} \texttt{section\_alpha\_locked.tex}, lines 21--27.

\textbf{Exact claim (line 24):}
``On the internal microsector $X = \mathbb{C}P^2 \times S^3$, we take a
Laplace-type operator given by the \emph{connection Laplacian}:''

\textbf{Equation (label: eq:connection-laplacian, lines 25--27):}
\[
  P = -g^{ij}\nabla_i\nabla_j
\]

\textbf{Clarification:} This is the \emph{connection Laplacian}
(also called the Bochner Laplacian or rough Laplacian), \emph{not}
the Dirac operator and \emph{not} the Hodge Laplacian.
It is a scalar-type operator (second order, no endomorphism term $E$
in the Gilkey sense beyond what is absorbed into the connection).

The paper explicitly distinguishes this from alternatives:
``Alternative operators would introduce additional terms proportional
to curvature derivatives, creating ppm-level ambiguities'' (line 34).
\end{tcolorbox}

%=============================================================================
\section{Question 2: Which Bundle?}
\label{sec:q2}
%=============================================================================

\begin{tcolorbox}[colback=yellow!5!white, colframe=yellow!50!black,
  title=Answer: LINE BUNDLE over CP2 with Kaehler curvature; TRIVIAL over S3]

\textbf{Source:} \texttt{section\_alpha\_locked.tex}, lines 30--31.

\textbf{Exact claim (lines 30--31):}
``The U(1) factor is implemented via twisting by a line bundle over
$\mathbb{C}P^2$ with curvature proportional to the K\"ahler form $\omega$,
taken trivial over $S^3$. This choice is minimal and convention-stable:
the K\"ahler form is parallel ($\nabla\omega = 0$), so derivative terms
in higher Seeley--DeWitt coefficients vanish automatically.''

\textbf{What is NOT stated:}
\begin{itemize}
\item The paper does \emph{not} specify which line bundle $\mathcal{O}(m)$
  the connection Laplacian acts on in this section.  Elsewhere in the paper,
  the twist bundle $E = \mathcal{O}(9) \oplus \mathcal{O}^{\oplus 5}$ appears
  in the Bridge Lemma context (for computing $k_{\max}$), but the gauge-kinetic
  operator bundle is described only as ``a line bundle over $\mathbb{C}P^2$
  with curvature proportional to $\omega$,'' which means $\mathcal{O}(1)$
  (up to normalization).
\item The paper does not state the rank of the bundle explicitly.
  ``Line bundle'' implies rank 1.
\end{itemize}
\end{tcolorbox}

%=============================================================================
\section{Question 3: Which Heat Kernel Coefficient?}
\label{sec:q3}
%=============================================================================

\begin{tcolorbox}[colback=yellow!5!white, colframe=yellow!50!black,
  title=Answer: $a_4$ (the fourth Seeley--DeWitt coefficient)]

\textbf{Source:} \texttt{section\_alpha\_locked.tex}, lines 15, 34, 53, 55.

The paper references $a_4$ at four locations:
\begin{enumerate}
\item Line 15: ``the predicted $\alpha$ must be stable without invoking
  subleading heat-kernel terms as ppm-level tuners.  Concretely, we choose
  a cutoff rule that prevents $a_6, a_8, \ldots$ from acting as free
  correction dials.''
\item Line 34: ``The gauge-kinetic extraction from $a_4$ is unambiguous
  with this operator choice.''
\item Line 53: ``the $a_6$ contribution would be $\sim 2\%$---far too
  large and requiring fine-tuned cancellation.''
\item Line 55: ``Preserves the leading $a_4$ gauge kinetic term.''
\end{enumerate}

\textbf{Convention:} The subscript numbering follows the convention where
$a_k$ corresponds to the coefficient of $t^{(k-n)/2}$ in the heat trace
asymptotic expansion in $n$ dimensions.  Here $a_4$ is the term that
produces the Yang--Mills action $\mathrm{Tr}(F^2)$ in dimension 4.
On $\mathbb{C}P^2 \times S^3$ (7-dimensional), $a_4$ enters through
the product formula / convolution of heat kernel coefficients.
\end{tcolorbox}

%=============================================================================
\section{Question 4: Standard Gilkey Formula or Modified?}
\label{sec:q4}
%=============================================================================

\begin{tcolorbox}[colback=yellow!5!white, colframe=yellow!50!black,
  title=Answer: IMPLICIT --- no explicit Gilkey formula is written out]

\textbf{Critical finding:} The paper \emph{never writes out} the
Gilkey formula for $a_4$.  It states that ``the gauge-kinetic extraction
from $a_4$ is unambiguous'' (line 34) but does not display the
curvature-invariant expression.

The paper does \emph{not} write anything like:
\[
a_4 = \frac{1}{360(4\pi)^{n/2}} \int \mathrm{tr}\bigl(
  c_1 R^2 + c_2 R_{ij}R^{ij} + c_3 R_{ijkl}R^{ijkl}
  + c_4 \Omega_{ij}\Omega^{ij} + \ldots\bigr)
\]

Instead, the paper's $\alpha$ derivation proceeds through an
\textbf{alternative route}: Chern--Simons partition function weights
on $S^3$ (appendix\_K.tex), combined with lattice Monte Carlo verification.
The spectral action language ($a_4$, plateau cutoff, Toeplitz truncation)
provides the \emph{framework} and \emph{justification} for the numerical
ingredients, but the actual numerical computation uses:
\begin{itemize}
\item The Chern--Simons weight function $w(k)$ (explicit)
\item The vacuum expectation $\langle k+2 \rangle$ (explicit)
\item The Toeplitz factors $(d-1)/d$ and $d^2/(d^2-1)$ (explicit)
\item The SM content parameters $N_{\rm species}$, $\mathrm{Tr}(Y^2)$, $g_F$ (tabulated)
\end{itemize}

Whether the Gilkey formula is \emph{standard} or \emph{modified} is
therefore not directly answerable from the text, because the Gilkey formula
is never explicitly invoked.
\end{tcolorbox}

%=============================================================================
\section{Question 5: Metric Normalization on CP2 and S3?}
\label{sec:q5}
%=============================================================================

\begin{tcolorbox}[colback=yellow!5!white, colframe=yellow!50!black,
  title=Answer: NOT EXPLICITLY STATED in section\_alpha\_locked.tex]

\texttt{section\_alpha\_locked.tex} does \emph{not} state the metric
normalization (e.g., whether $\mathbb{C}P^2$ uses Fubini--Study with
sectional curvature in $[1,4]$, or $S^3$ has radius $R = 1$ or
some other value).

\texttt{appendix\_K.tex} (line 39--42) gives the heat kernel on $S^3$
with radius $R$:
\[
K(t; S^3) = \sum_{n=0}^{\infty} (n+1)^2 e^{-n(n+2)t/R^2}
\]
but does not fix $R$ to a particular value.

The standalone PDF similarly does not state metric normalizations
explicitly.  The curvature invariants of $\mathbb{C}P^2$ and $S^3$
are never written out in any of the three source documents.
\end{tcolorbox}

%=============================================================================
\section{Question 6: Intermediate Numerical Values for $a_4$?}
\label{sec:q6}
%=============================================================================

\begin{tcolorbox}[colback=yellow!5!white, colframe=yellow!50!black,
  title=Answer: NO explicit $a_4$ numerical value is given]

The paper does \emph{not} state a numerical value for $a_4$ before
applying SM traces.  No intermediate value like ``$a_4 = X$ on
$\mathbb{C}P^2 \times S^3$'' appears anywhere.

What \emph{is} given numerically:
\begin{itemize}
\item $\beta_{U(1)} = \langle k+2 \rangle_{k_{\max}=60} = 3.7969 \approx 3.80$
  (appendix\_K.tex, line 78; PDF Eq.~15)
\item $\Lambda^3 = k \cdot (k+3)/(k+4) = 60 \cdot 63/64 = 885.9375$
  (section\_alpha\_locked.tex, line 79; also line 210)
\item $(d-1)/d = 63/64$ (section\_alpha\_locked.tex, line 196)
\item $\varepsilon_{\rm adj} = 4096/4095 = 1.000244\ldots$
  (section\_alpha\_locked.tex, line 119)
\item $w = N_{\rm species}/(g_F \cdot \mathrm{Tr}(Y^2)) = 7/80$
  (section\_alpha\_locked.tex, line 200)
\item Final result: $\alpha^{-1} = 137.03599985$
  (section\_alpha\_locked.tex, line 145, 203, 253)
\end{itemize}

The derivation chain table (lines 184--207) lists all components
but never isolates an intermediate ``$a_4$ value.''
\end{tcolorbox}

%=============================================================================
\section{Question 7: References to Gilkey, Vassilevich, or Chamseddine--Connes?}
\label{sec:q7}
%=============================================================================

\begin{tcolorbox}[colback=yellow!5!white, colframe=yellow!50!black,
  title=Answer: NONE in section\_alpha\_locked.tex or appendix\_K.tex]

\texttt{section\_alpha\_locked.tex}: Contains zero bibliographic references
to Gilkey, Vassilevich, or Chamseddine--Connes.  The only references are
internal cross-references (Sec.~, App.~).

\texttt{appendix\_K.tex}: Contains zero references to Gilkey, Vassilevich,
or Chamseddine--Connes.  References Witten 1989 [6] for the CS partition
function, Griffiths--Harris [7] for algebraic geometry, and
Atiyah--Singer [8] for the index theorem.

The standalone PDF references [1]--[10], none of which are Gilkey,
Vassilevich, or Chamseddine--Connes.  The references are: Feynman [1],
Berry [2], Wilczek--Zee [3], Wilson [4], Creutz [5], Witten [6],
Griffiths--Harris [7], Atiyah--Singer [8], Alcock (DFD book) [9],
Alcock (simulation code) [10].

\textbf{Note:} A supplementary file
\texttt{DFD\_Standard\_Model\_pasted\_note.tex} (line 441) does reference
``Gilkey 1975'' in its bibliography, and \texttt{a4\_SeeleyDeWitt.tex}
(a research output file, not part of the paper itself) discusses the
Gilkey formula.  But these are \emph{not} in the published/submitted paper.
\end{tcolorbox}

%=============================================================================
\section{Question 8: What Does ``the gauge-kinetic extraction from $a_4$
  is unambiguous'' Mean?}
\label{sec:q8}
%=============================================================================

\begin{tcolorbox}[colback=yellow!5!white, colframe=yellow!50!black,
  title=Answer: It means the connection Laplacian with parallel curvature
  produces a unique $\mathrm{Tr}(F^2)$ term]

\textbf{Source:} \texttt{section\_alpha\_locked.tex}, lines 33--34.

\textbf{Full context (lines 33--34):}
``Why this is locked. The gauge-kinetic extraction from $a_4$ is unambiguous
with this operator choice.  Alternative operators would introduce additional
terms proportional to curvature derivatives, creating ppm-level ambiguities.
The connection Laplacian with parallel curvature eliminates this freedom.''

\textbf{Interpretation:}
The claim is that for the \emph{connection Laplacian} (as opposed to, say,
a Dirac-type operator or a Laplacian with additional endomorphism terms),
the $a_4$ coefficient contains a $\mathrm{Tr}(\Omega_{ij}\Omega^{ij})$ term
(where $\Omega$ is the bundle curvature) that directly yields the
gauge-kinetic $\mathrm{Tr}(F^2)$ action, and because the K\"ahler form is
parallel ($\nabla\omega = 0$), there are no additional derivative-of-curvature
terms that could modify this extraction at the ppm level.

\textbf{No specific formula is referenced.}  The paper does not point to
a specific equation number in Gilkey, Vassilevich, or any external source.
It is a qualitative statement about why the operator choice is unambiguous,
not a quantitative formula.
\end{tcolorbox}

%=============================================================================
\section{Question 9: Closed-Form or Numerical Result?}
\label{sec:q9}
%=============================================================================

\begin{tcolorbox}[colback=yellow!5!white, colframe=yellow!50!black,
  title=Answer: HYBRID --- closed-form structure with one numerically
  computed input]

The derivation has two distinct layers:

\textbf{Layer 1: Closed-form algebraic structure.}
All of the following are exact rational or algebraic expressions:
\begin{itemize}
\item $d = k + 4 = 64$ (exact integer)
\item $(d-1)/d = 63/64$ (exact rational)
\item $d^2/(d^2-1) = 4096/4095$ (exact rational)
\item $\Lambda^3 = k(k+3)/(k+4) = 885.9375$ (exact rational)
\item $N_{\rm species} = 7$, $\mathrm{Tr}(Y^2) = 10$, $g_F = 8$ (exact integers)
\item $w = 7/80$ (exact rational)
\end{itemize}

\textbf{Layer 2: One numerically computed input.}
\[
\beta_{U(1)} = \langle k+2 \rangle_{k_{\max}=60}
= \frac{\sum_{k=0}^{59}(k+2)\,w(k)}{\sum_{k=0}^{59} w(k)} = 3.7969\ldots
\]
where $w(k) = \frac{2}{k+2}\sin^2\!\bigl(\frac{\pi}{k+2}\bigr)$.
This is a finite sum of trigonometric terms---computable to arbitrary
precision, but not a simple closed form.

\textbf{Final result:}
The sub-ppm value $\alpha^{-1} = 137.03599985$ (section\_alpha\_locked.tex,
line 145) is stated as exact given the inputs.  The standalone PDF gives
a more conservative $\alpha^{-1} = 137.036 \pm 0.5$ (appendix\_K.tex,
line 95), reflecting lattice simulation uncertainty.

\textbf{The lattice Monte Carlo results} (standalone PDF, 86 runs) are
purely numerical, providing verification at the $\sim 1\%$ level.
\end{tcolorbox}

%=============================================================================
\section{Question 10: External Computation Tools or Code?}
\label{sec:q10}
%=============================================================================

\begin{tcolorbox}[colback=yellow!5!white, colframe=yellow!50!black,
  title=Answer: YES --- Zenodo code repository explicitly referenced]

\textbf{Source:} Standalone PDF, page 3 and page 23.

\textbf{Code and data:} \url{https://doi.org/10.5281/zenodo.19173548}
(Reference [10] in the PDF).

The PDF includes two inline code listings:
\begin{enumerate}
\item \textbf{Listing 1 (page 23):} Python code computing
  $\langle k+2 \rangle$ for various $k_{\max}$ values.  Uses the weight
  function \texttt{w(k) = (2.0/(k+2)) * (math.sin(math.pi/(k+2)))**2}.
\item \textbf{Listing 2 (page 24):} Python function
  \texttt{alpha\_wilson(ku, ks)} computing $\alpha_W$ from measured
  stiffnesses.  Uses \texttt{g1 = 1.0/ku} (U(1) standard) and
  \texttt{g2 = 4.0/ks} (SU(2) Wilson normalization).
\end{enumerate}

\textbf{No symbolic algebra system} (Mathematica, Maple, etc.) is
referenced.  The computations are simple enough for hand calculation
or short Python scripts.
\end{tcolorbox}

%=============================================================================
\section{Summary of Critical Gaps for Reproduction}
\label{sec:gaps}
%=============================================================================

The following items are \emph{not} stated explicitly in any of the three
source documents and would need to be established by an agent attempting
to reproduce the computation:

\begin{longtable}{p{0.35\textwidth}p{0.55\textwidth}}
\toprule
\textbf{Missing Item} & \textbf{Status in Paper} \\
\midrule
Explicit Gilkey formula for $a_4$ &
  Never written out.  Referred to qualitatively. \\
Curvature invariants of $\mathbb{C}P^2$ &
  Never stated ($R$, $R_{ij}R^{ij}$, $R_{ijkl}R^{ijkl}$, Weyl). \\
Curvature invariants of $S^3$ &
  Never stated. \\
Metric normalization of $\mathbb{C}P^2$ &
  Not specified (Fubini--Study assumed but scale not given). \\
Metric normalization of $S^3$ &
  Radius $R$ appears symbolically but is not fixed. \\
Product formula for $a_4$ on $M_1 \times M_2$ &
  Never stated (the convolution $a_4 = \sum a_j \otimes a_{4-j}$ is implicit). \\
How $\beta_{U(1)} = 3.80$ connects to $\kappa_{U(1)}$ &
  The standalone PDF uses lattice measurement; section\_alpha\_locked.tex
  uses the spectral action framework with the derivation chain table,
  but the bridging formula $\kappa_{U(1)} = f(\beta_{U(1)}, \ldots)$ is
  not written as a single equation. \\
The role of $N_{\rm species} = 7$ &
  Stated as ``SM SU(2) components'' in the derivation chain table (line 197)
  but not derived or explained further in this section. \\
The role of $\mathrm{Tr}(Y^2) = 10$ &
  Stated as ``SM hypercharges'' (line 198).  Note: appendix\_Z.tex gives
  $\mathrm{Tr}(Y^2) = 10/3$ per generation (standard SM computation),
  so the value 10 apparently includes a factor of 3 for $N_{\rm gen}$. \\
The role of $g_F = 8$ &
  Stated as coming from ``Spectral triple ($J \times \gamma \times
  \mathbb{C}$)'' (line 199).  Not derived in any of the three source
  documents. \\
\bottomrule
\end{longtable}

%=============================================================================
\section{The Two Derivation Routes}
\label{sec:two-routes}
%=============================================================================

A critical observation is that the DFD corpus contains \textbf{two distinct
derivation routes} for $\alpha$, which must not be conflated:

\subsection{Route 1: Spectral Action / Toeplitz (section\_alpha\_locked.tex)}

This is the ``convention-locked'' derivation:
\begin{enumerate}
\item Operator: connection Laplacian on $\mathbb{C}P^2 \times S^3$
\item Regulator: plateau cutoff (kills $a_6, a_8, \ldots$)
\item Truncation: Toeplitz at $d = 64$, giving $(d-1)/d = 63/64$
\item Fork: regular-module $\mathcal{H}_F = M_d(\mathbb{C})$, giving
  boost $d^2/(d^2-1) = 4096/4095$
\item SM content: $N_{\rm species} = 7$, $\mathrm{Tr}(Y^2) = 10$, $g_F = 8$
\item Hypercharge weight: $w = 7/80$
\item Result: $\alpha^{-1} = 137.03599985$ (sub-ppm)
\end{enumerate}

This route claims sub-ppm precision but never writes out the explicit
$a_4$ formula.

\subsection{Route 2: Chern--Simons + Lattice MC (standalone PDF / appendix\_K.tex)}

This is the lattice verification route:
\begin{enumerate}
\item UV cutoff: $k_{\max} = \chi(\mathbb{C}P^2, E) = 60$ (Bridge Lemma)
\item Chern--Simons vacuum: $\beta_{U(1)} = \langle k+2 \rangle = 3.80$
\item Lattice ratio: $\beta_{SU(2)}/\beta_{U(1)} = (n_2/n_1) \times N_{\rm gen} = 6$
\item Stiffness ratio: $\kappa_{U(1)}/\kappa_{SU(2)} = 1/2$
\item Lattice MC at $(\beta_{U(1)}, \beta_{SU(2)}) = (3.80, 22.80)$
\item Stiffness extraction: $\kappa = F''(0)/V$
\item Coupling extraction: $g_1^2 = 1/\kappa_{U(1)}$,
  $g_2^2 = 4/\kappa_{SU(2)}$ (Wilson normalization)
\item $\alpha_W$ from electroweak mixing formula
\item Result: $\alpha \approx 1/137$ within $\sim 1\%$
\end{enumerate}

This route has $\sim 1\%$ precision from lattice statistics.

\subsection{Relationship}

The standalone PDF (Section VI.C, page 21) explicitly states:
``The parent theory [9] contains additional machinery beyond what is
presented here.  In particular, the full derivation uses a
Toeplitz-truncated spectral action on $\mathbb{C}P^2 \times S^3$ and
resolves a forced binary fork between a regular-module microsector and
a fermion-representation microsector via a no-hidden-knobs policy.
Only the regular-module branch survives, yielding
$\alpha^{-1} = 137.03599985$ at sub-ppm precision.''

%=============================================================================
\section{Complete Equation Inventory}
\label{sec:equations}
%=============================================================================

\subsection{section\_alpha\_locked.tex}

\begin{longtable}{p{0.25\textwidth}p{0.15\textwidth}p{0.50\textwidth}}
\toprule
\textbf{Label} & \textbf{Line} & \textbf{Content} \\
\midrule
eq:connection-laplacian & 25--27 & $P = -g^{ij}\nabla_i\nabla_j$ \\
eq:spectral-action-plateau & 41--43 & $S = \mathrm{Tr}\,f(P/\Lambda^2)$ \\
(unnumbered) & 48--50 & $f_{-2} = f_{-4} = \cdots = 0$ \\
eq:d-dimension & 65--67 & $d = \dim H^0(\mathbb{C}P^1, \mathcal{O}(m)) = k + 4$ \\
(unnumbered) & 71--73 & $K_{\mathbb{C}P^2} = \mathcal{O}(-3)$ \\
eq:lambda-cutoff & 78--80 & $\Lambda^3 = k \cdot (k+3)/(k+4)$ \\
eq:HF-regular & 93--95 & $\mathcal{H}_F := A = M_d(\mathbb{C})$ \\
eq:tr-dem & 103--105 & $\mathrm{tr}_{\rm dem}(\cdot) = (1/d^2)\mathrm{Tr}$ \\
(unnumbered) & 109--111 & $\mathrm{tr}_{\mathfrak{su}}(\cdot) = (1/(d^2-1))\mathrm{Tr}_{\mathfrak{su}}$ \\
eq:boost-factor & 113--115 & $\varepsilon_{\rm adj}^{(A)} = d^2/(d^2-1)$ \\
(unnumbered) & 118--120 & $\varepsilon_{\rm adj}^{(A)} = 4096/4095$ \\
eq:drop-factor & 125--127 & $\varepsilon_{\rm adj}^{(B)} = (d^2-1)/d^2 = 4095/4096$ \\
\bottomrule
\end{longtable}

\subsection{Standalone PDF Key Equations}

\begin{longtable}{p{0.10\textwidth}p{0.15\textwidth}p{0.65\textwidth}}
\toprule
\textbf{Eq.} & \textbf{Page} & \textbf{Content} \\
\midrule
(1) & 5 & $\alpha = e^2/(4\pi\epsilon_0\hbar c) \approx 1/137.036$ \\
(2) & 6 & $\mathbf{a} = (c^2/2)\nabla\psi$ \\
(3) & 6 & $Z_{SU(2)_k}(S^3) = \sqrt{2/(k+2)}\sin(\pi/(k+2))$ \\
(5) & 6 & $w(k) = |Z|^2 = (2/(k+2))\sin^2(\pi/(k+2))$ \\
(7) & 7 & $\mathcal{L}_{\rm stiff}^{(r)} = -(\kappa_r/2)\mathrm{Tr}\,F_{ij}^{(r)}F_{ij}^{(r)}$ \\
(8) & 7 & $\kappa_r = n_r\kappa_0$ \\
(9) & 7 & $\kappa_{U(1)}/\kappa_{SU(2)} = n_1/n_2 = 1/2$ \\
(10) & 7 & $1/e^2 = 1/g_1^2 + 1/g_2^2$ \\
(12) & 8 & $g_1^2 = 1/\kappa_{U(1)}$, $g_2^2 = 4/\kappa_{SU(2)}$ \\
(13) & 8 & $\alpha_W = (1/\kappa_{U(1)})(4/\kappa_{SU(2)})/((1/\kappa_{U(1)}) + (4/\kappa_{SU(2)})) \cdot 1/(4\pi)$ \\
(14) & 8 & $k_{\max} = \chi(\mathbb{C}P^2, \mathcal{O}(9)\oplus\mathcal{O}^{\oplus 5}) = \binom{11}{2} + 5 = 60$ \\
(15) & 8 & $\langle k_{\rm eff}\rangle = \sum_{k=0}^{59}(k+2)w(k) / \sum w(k) = 3.7969 \approx 3.80$ \\
(16) & 9 & $\beta_{U(1)} = \langle k_{\rm eff}\rangle_{k_{\max}=60} = 3.80$ \\
(17) & 10 & $\beta_{SU(2)}/\beta_{U(1)} = (n_2/n_1) \times N_{\rm gen} = 2 \times 3 = 6$ \\
(19) & 13 & Lattice action $S$ (full expression) \\
(20) & 14 & $\kappa = F''(0)/V$ (stiffness extraction) \\
\bottomrule
\end{longtable}

%=============================================================================
\section{Key Numerical Values Stated in the Papers}
\label{sec:numerics}
%=============================================================================

\begin{longtable}{p{0.30\textwidth}p{0.20\textwidth}p{0.40\textwidth}}
\toprule
\textbf{Quantity} & \textbf{Value} & \textbf{Source Location} \\
\midrule
$k_{\max}$ & 60 & PDF Eq.~14; appendix\_K line 68 \\
$d = k + 4$ & 64 & alpha\_locked line 66 \\
$(d-1)/d$ & $63/64$ & alpha\_locked line 196 \\
$d^2/(d^2-1)$ & $4096/4095$ & alpha\_locked line 119 \\
$\Lambda^3$ & $885.9375$ & alpha\_locked line 210 \\
$\beta_{U(1)}$ & $3.7969 \approx 3.80$ & appendix\_K line 78; PDF Eq.~15 \\
$\beta_{SU(2)}$ & $22.80$ & PDF Eq.~18 \\
$\beta_{SU(2)}/\beta_{U(1)}$ & $6$ & PDF Eq.~17 \\
$n_2/n_1$ & $2$ & PDF Eq.~9 \\
$N_{\rm gen}$ & $3$ & PDF Sec.~III.C \\
$\kappa_{U(1)}/\kappa_{SU(2)}$ & $0.5$ (predicted) & PDF Eq.~9 \\
$\kappa_{U(1)}/\kappa_{SU(2)}$ & $0.495 \pm 0.020$ (measured) & PDF Sec.~V.E \\
$N_{\rm species}$ & $7$ & alpha\_locked line 197 \\
$\mathrm{Tr}(Y^2)$ & $10$ & alpha\_locked line 198 \\
$g_F$ & $8$ & alpha\_locked line 199 \\
$w = N_{\rm sp}/(g_F \cdot \mathrm{Tr}(Y^2))$ & $7/80$ & alpha\_locked line 200 \\
$\alpha^{-1}$ (Route 1) & $137.03599985$ & alpha\_locked line 145 \\
$\alpha^{-1}$ (Route 2) & $137.036 \pm 0.5$ & appendix\_K line 95 \\
$\alpha_W$ at $L=6$, $\beta=3.80$ & $0.007297$ & PDF Table VI \\
$\alpha_W$ at $L=12$, $\beta=3.80$ & $0.007291$ & PDF Table VI \\
$\alpha_{\rm phys}$ & $0.0072973525693$ & (experimental) \\
\bottomrule
\end{longtable}

\end{document}
